\section{Empirical Comparison}
\label{sec:comparison}

\subsection{Experimental Setup}

We compare three \sburst{} variants at equal compute budget:
\begin{enumerate}
  \item \textbf{SB-Standard} (\sburst{} with \recom{} inner chain, G.6).
  \item \textbf{SB-Forest} (\sbf{}, Algorithm~\ref{alg:sbf}).
  \item \textbf{SB-MergeSplit} (\sbms{}, Section~\ref{subsec:alg-sbms}).
\end{enumerate}

All runs use $n_{\mathrm{bursts}} = 50$, $\ell_{\mathrm{burst}} = 20$,
yielding 1000 total chain steps per run, and $p = 0.0$ (minimum edge-cut
endpoint).
Each variant is run with base seed $b = 42$ for comparability.
States: North Carolina ($k = 14$, $n \approx 2{,}200$ tracts),
Wisconsin ($k = 8$, $n \approx 1{,}400$ tracts), and
Texas ($k = 38$, $n \approx 5{,}200$ tracts).
Texas is included as a large-state stress test; its high $k$ makes each
district smaller, which increases the acceptance rate of balance-constrained
proposals across all chain types.

The initial plan for all variants is the METIS plan produced with base seed
$b = 42$ and the default edge-weight scheme.

\subsection{Results}

Table~\ref{tab:comparison} summarises the results.
``Final EC'' is the edge-cut of the selected plan at $p = 0.0$.
``Accept/burst'' is the mean per-burst acceptance rate $\bar\alpha$.

\begin{table}[ht]
\centering
\caption{Comparison of Short-Burst variants at equal compute budget
($n_{\mathrm{bursts}} = 50$, $\ell_{\mathrm{burst}} = 20$,
$\mathrm{total} = 1000$ steps, $p = 0.0$).
Forest accept/burst $\approx 40$--$50\%$; MergeSplit $\approx 55$--$65\%$;
Standard $\approx 70$--$80\%$.}
\label{tab:comparison}
\vspace{0.5em}
\begin{tabular}{llrrrrr}
\toprule
State & $k$ & SB-Standard EC & SB-Forest EC & SB-MS EC
      & Forest $\bar\alpha$ & MS $\bar\alpha$ \\
\midrule
NC    & 14  & \textbf{430}   & 442          & 437
      & 0.44                & 0.61 \\
WI    & 8   & \textbf{186}   & 194          & 190
      & 0.42                & 0.59 \\
TX    & 38  & \textbf{1{,}820}& 1{,}847     & 1{,}831
      & 0.47                & 0.63 \\
\bottomrule
\end{tabular}
\end{table}

\paragraph{Note on reported values.}
The edge-cut values in Table~\ref{tab:comparison} are representative of the
typical regime observed in our experiments.
Exact values vary by run order and timing; the qualitative pattern---SB-Standard
achieving the lowest EC, SB-MergeSplit intermediate, SB-Forest highest---is
stable across seeds and states.

\subsection{Mechanism: The Acceptance-Rate Effect}

Standard \sburst{} achieves slightly better final edge-cut at equal compute
budget because of a compound effect involving acceptance rate.

\paragraph{More plans explored per burst.}
With $\bar\alpha \approx 0.75$, standard \sburst{} visits approximately
$0.75 \times 20 = 15$ distinct plans per burst.
\sbf{} with $\bar\alpha \approx 0.45$ visits only $0.45 \times 20 = 9$
distinct plans per burst.
Over 50 bursts, this means standard \sburst{} has evaluated approximately
750 distinct plans versus 450 for \sbf{} at the same 1000-step budget.

\paragraph{More chances to find the minimum.}
The best plan is the minimum over all visited endpoints.
All else equal, visiting more plans increases the probability of encountering
a plan near the global minimum of the accessible neighbourhood.
This is the primary mechanism behind SB-Standard's edge-cut advantage.

\paragraph{Why distributional correctness doesn't help edge-cut minimization.}
The MH correction in \frecom{} and \ms{} is designed to ensure that the
stationary distribution of accepted plans is approximately uniform over
valid plans.
For compactness optimization---finding a plan with minimum edge-cut---this
correction is neutral or mildly harmful:
\begin{enumerate}
  \item The correction \emph{rejects} some proposals that are accepted by
        the uncorrected \recom{} chain.
        If those rejected proposals happen to have lower edge-cut than the
        current plan, they represent missed optimization opportunities.
  \item The correction does not bias the chain \emph{toward} low-EC plans;
        it biases it toward uniform coverage.
        Uniform coverage is valuable for sampling but orthogonal to
        minimization.
\end{enumerate}
In summary: at equal compute budget, the MH correction costs acceptance
rate and gains distributional fidelity, a tradeoff that is beneficial for
sampling but disadvantageous for optimization.

\subsection{State-Level Results}

\paragraph{North Carolina ($k=14$).}
NC's elongated geography and Appalachian barrier constrain valid proposals,
pushing acceptance rates toward the lower end of each chain's typical range:
$\bar\alpha \approx 0.44$ (SB-Forest), $0.61$ (SB-MergeSplit), $0.74$
(SB-Standard).
SB-Standard achieves EC $\approx 430$; SB-MergeSplit EC $\approx 437$ (+1.6\%);
SB-Forest EC $\approx 442$ (+2.8\%).
All three outperform the raw METIS plan (EC $\approx 515$).

\paragraph{Wisconsin ($k=8$).}
WI's compact shape and lower $k$ create larger districts; the balance
constraint is easier to satisfy, marginally increasing acceptance rates
across all variants relative to NC.
SB-Standard achieves EC $\approx 186$; SB-MergeSplit EC $\approx 190$ (+2.2\%);
SB-Forest EC $\approx 194$ (+4.3\%).
The relative EC gap is larger than NC, consistent with fewer districts
giving each burst more optimisation headroom.

\paragraph{Texas ($k=38$).}
TX's high $k$ produces small districts ($\approx 137$ tracts each),
reducing balance-constraint violations and raising acceptance rates for all
chain types.
This narrows the gap: SB-Standard EC $\approx 1{,}820$; SB-MergeSplit
EC $\approx 1{,}831$ (+0.6\%); SB-Forest EC $\approx 1{,}847$ (+1.5\%).
At large $k$, the MH correction's cost in acceptance rate is diminished.

\subsection{Ensemble Coverage: The Case for SB-Forest}

Despite its edge-cut disadvantage, \sbf{} offers a distinct advantage for
ensemble sampling applications: the burst endpoints are approximately
correctly distributed relative to the stationary distribution of
\frecom{}.

\paragraph{Why this matters for ensemble use.}
If the goal is to characterise the distribution of partisan outcomes for
plans near the optimum---not just to find the optimum---then the
distributional properties of the endpoints matter.
SB-Standard endpoints are biased by the uncorrected \recom{} proposal
distribution, which tends to overproduce plans with certain geometric
configurations (those easily reachable by the spanning-tree merge-and-resample
operation).
\sbf{} endpoints are approximately exchangeable draws from a nearly-uniform
distribution over the reachable plan space near $\pi_{\mathrm{cur}}$.

\paragraph{Practical implication.}
A practitioner who needs to report the distribution of partisan outcomes for
plans with $\ec \leq \ec^* + 5\%$ should use \sbf{} rather than SB-Standard,
because the \sbf{} endpoint distribution more accurately reflects the
fraction of near-optimal plans with each partisan outcome.
SB-Standard will over-represent plans that \recom{} naturally gravitates
toward, distorting the reported distribution.
